\section{Continuous Integration}

Vediamo infine come ho impostato la continuous integration sul progetto. L'idea è che,
ogni volta che si fa un push o si apre una pull request, i test partano da soli, così se
qualcuno rompe qualcosa ce ne accorgiamo subito. Il tutto sta nel workflow
\texttt{.github/workflows/ci.yml}: fa la build Maven completa (lasciando fuori i test
taggati \texttt{ui}, che hanno bisogno di Chrome e dell'app avviata) e alla fine salva
il report JaCoCo come artifact.

\begin{lstlisting}[language=bash, caption={.github/workflows/ci.yml (estratto)}]
on:
  push:
  pull_request:

jobs:
  build-and-test:
    runs-on: ubuntu-latest
    steps:
      - uses: actions/checkout@v4
      - uses: actions/setup-java@v4
        with:
          distribution: temurin
          java-version: '17'
      - run: mvn -B -f java/pom.xml verify
      - uses: actions/upload-artifact@v4
        with:
          name: jacoco-report
          path: java/target/site/jacoco/
\end{lstlisting}

\screenshot{ci_success}{Esecuzione CI riuscita}{%
Dopo aver creato il repository GitHub e fatto il primo push, andare sulla tab
\textbf{Actions} del repository. Screenshot di un run completato con successo
(spunta verde), con tutti gli step espansi visibili (checkout, setup JDK, build+test).}

\screenshot{ci_failure}{Esecuzione CI fallita (dimostrazione)}{%
Per dimostrare che la CI intercetta davvero le regressioni: modificare
temporaneamente un metodo del core (es. far restituire sempre \texttt{TipoPosto.NESSUNO}
da \texttt{assegnaPosto}), fare commit/push su un branch di prova, e catturare lo
screenshot del run \textbf{fallito} in rosso con il log dei test JUnit falliti
espanso. Poi ripristinare il codice corretto con un nuovo commit e, se vuoi, uno
screenshot finale che torna verde.}
